%% -*- coding: utf-8 -*-

\section{诗人简介}

\newcounter{myterm}
\begin{description}[
  before={\setcounter{myterm}{0}}, % 关键：每次进入这个列表时，把序号清零
  font={\stepcounter{myterm}\arabic{myterm}.~\bfseries}, % 关键：每次遇到 \item，数字加1并输出，然后保持名词粗体
  labelsep=0.5em,
  leftmargin=3em % 稍微增加一点左侧缩进，给序号留出空间
]
  \item [钱仲联] （1908年—2003年），初名萼孙，字仲联，号梦苕，以字行，江苏常熟人，原籍浙江省吴兴县，中国文学研究家，任教于大夏大学、无锡国学专修学校、南京中央大学、苏州大学等校。有“昆仑万象一代诗豪”、“清诗功臣”之美誉。
  \item [刘世南] （1923年10月-2021年8月1日），江西吉安人，文史学者、诗人。曾任江西师范大学文学院副教授、硕士生导师，1989年退休，2006年补评为教授，任《豫章丛书》点校首席学术顾问及江西省古籍整理中心组成员。\par
        幼年随父研习古籍，高一辍学后长期在吉安、新建等地任中学语文教师。1979年因发表《谈古文的标点、注释和翻译》等论文引发学界关注，同年调入江西师范大学任教。代表专著《清诗流派史》由台北文津出版社和人民文学出版社先后出版，被学界视为清诗研究经典著作，另著有《清诗三百首详注》《黄遵宪诗选注》《清文选》《谷梁传直解》《古文观止译注》《在学术殿堂外》《大螺居诗存》《大螺居诗文存》《师友偶记》等。旧体诗作入选《海岳弦歌集》等选集，并为高校撰拟楹联铭文。2021年8月1日在南昌逝世，享年99岁。
  \item [\hypertarget{XingFang}{邢昉}] （1590年—1653年），字孟贞，一字石湖。江苏高淳人。\par

        曾祖邢符，祖邢尚宾，父邢一淳，科场俱不得意。明末诸生，屡试不中，第六次参加乡试，主考官称其文笔太狂﹐愤而作《太狂篇》，从此绝意仕途。与方文友好，两人有诗往来。崇祯十年丁丑（1637年），应杨文骢之请赴华亭幕府。崇祯十三年庚辰（1640年），邢昉、方文等人于金陵结社。崇祯十四年辛巳（1641年），方文赴庐州投靠蔡如蘅，邢昉有诗相送。崇祯十六年癸未（1643年），邢昉旅居宣城。崇祯十七年甲申（1644年）八月，母赵孺人逝。\par
        邢昉是明代著名诗人，终身布衣，为诗清真古澹，以《故宫燕》最著名，王士禛在《渔洋诗话》独推邢昉“韦、柳门庭人”及“布衣诗人第一”。吴敬梓曾拜访其石臼湖畔故居，谒邢昉墓。著有《宛游草》、《石臼集》。
\end{description}


%%% Local Variables:
%%% TeX-engine: xetex
%%% TeX-master: "Qing_Poem_300_Compare"
%%% End:
